% =====================================================================
%  BAO CAO KIEM THU HIEU NANG AI (LOAD TEST voi k6)
%  Bien dich bang: xelatex bao_cao_hieu_nang_ai.tex
%  (XeLaTeX + fontspec ho tro tieng Viet UTF-8 san, khong can vntex)
% =====================================================================
\documentclass[12pt,a4paper]{article}

\usepackage{fontspec}          % can XeLaTeX; font Latin Modern mac dinh ho tro tieng Viet
\usepackage[a4paper,margin=2.5cm]{geometry}
\usepackage{booktabs}
\usepackage{array}
\usepackage{longtable}
\usepackage{xcolor}
\usepackage{colortbl}
\usepackage{hyperref}
\hypersetup{colorlinks=true, linkcolor=blue!60!black, urlcolor=blue!60!black}

\definecolor{ok}{RGB}{0,140,60}
\definecolor{warn}{RGB}{200,140,0}
\definecolor{bad}{RGB}{200,40,40}
\definecolor{head}{RGB}{30,60,110}

\newcommand{\ok}[1]{\textcolor{ok}{#1}}
\newcommand{\warn}[1]{\textcolor{warn}{#1}}
\newcommand{\bad}[1]{\textcolor{bad}{#1}}

\title{\textbf{Báo cáo Kiểm thử Hiệu năng Hệ thống AI Text-to-SQL}\\[4pt]
\large Load Test với k6 --- Ứng dụng Vanna (Ollama \texttt{llama3.2} + PostgreSQL Neon)}
\author{Đức Anh \\ \small \texttt{ducanhtrungvan@gmail.com}}
\date{Ngày 04 tháng 08 năm 2026}

\begin{document}
\maketitle

\begin{abstract}
\noindent Báo cáo trình bày kết quả kiểm thử hiệu năng (load test) của hệ thống hỏi đáp
dữ liệu bằng ngôn ngữ tự nhiên (Text-to-SQL) xây dựng trên nền Vanna, sử dụng mô hình
ngôn ngữ lớn Ollama \texttt{llama3.2} chạy cục bộ và cơ sở dữ liệu PostgreSQL (Neon).
Công cụ \textbf{k6} được dùng để mô phỏng số lượng người dùng đồng thời tăng dần
(5, 15, 30 và 50 người dùng). Kết quả cho thấy hệ thống hoạt động ổn định dưới
\textbf{10 người dùng đồng thời}, nhưng \textbf{bão hoà} khi vượt ngưỡng này: thông lượng
(throughput) không tăng thêm mà độ trễ và tỉ lệ lỗi tăng phi tuyến. Ở mức 50 người dùng,
độ trễ p95 đạt \textbf{139 giây} và tỉ lệ lỗi lên tới \textbf{13\%}.
\end{abstract}

\section{Mục tiêu}
Đánh giá khả năng chịu tải của endpoint hỏi đáp AI khi số lượng người dùng đồng thời
tăng cao, xác định \emph{điểm bão hoà} và đề xuất hướng cải thiện.

\section{Môi trường và phương pháp}
\begin{table}[h]
\centering
\begin{tabular}{@{}>{\bfseries}l l@{}}
\toprule
Hạng mục & Chi tiết \\
\midrule
Công cụ kiểm thử & k6 v2.1.0 (Windows/amd64) \\
Endpoint & \texttt{POST /api/vanna/v2/chat\_sse} (streaming SSE) \\
Mô hình AI & Ollama \texttt{llama3.2} (chạy cục bộ) \\
Cơ sở dữ liệu & PostgreSQL --- Neon (\texttt{qlsp\_backup}) \\
Máy chủ ứng dụng & Uvicorn (1 worker), FastAPI \\
Loại truy vấn & Câu hỏi tiếng Việt (đếm / tổng hợp) \\
Tiêu chí thành công & HTTP 200 và stream kết thúc bằng \texttt{[DONE]} \\
Timeout mỗi request & 180 giây \\
Baseline (1 người dùng) & $\approx 2{,}5$ giây/truy vấn \\
\bottomrule
\end{tabular}
\caption{Cấu hình môi trường kiểm thử.}
\end{table}

\noindent\textbf{Lưu ý kỹ thuật:} Mỗi câu hỏi kích hoạt \textbf{2 lần gọi mô hình ngôn ngữ}
(một lần để sinh câu lệnh SQL, một lần để diễn giải kết quả). Mỗi request dùng một
\texttt{conversation\_id} riêng biệt nhằm tránh cơ chế khoá theo hội thoại (HTTP 409) của máy chủ.

\section{Kịch bản kiểm thử}
Hai kịch bản tăng tải (\texttt{ramping-vus}) được thực hiện:
\begin{itemize}
  \item \textbf{Kịch bản A --- Tăng dần:} $5 \to 15 \to 30$ người dùng đồng thời trong $\approx 3$ phút 35 giây.
  \item \textbf{Kịch bản B --- Cao điểm 50:} tăng lên $50$ người dùng và giữ trong 120 giây.
\end{itemize}

\section{Kết quả tổng hợp}

\begin{table}[h]
\centering
\renewcommand{\arraystretch}{1.25}
\begin{tabular}{@{}l c c c@{}}
\toprule
\rowcolor{head!12}
\textbf{Chỉ số} & \textbf{1 user (baseline)} & \textbf{30 user (A)} & \textbf{50 user (B)} \\
\midrule
Tổng số request      & 1     & 177   & 132 \\
Thông lượng (req/giây) & 0,40 & \textbf{0,75} & \textbf{0,70} \\
Tỉ lệ thành công     & 100\% & \ok{95,0\%} & \warn{87,0\%} \\
Số lỗi (rớt kết nối / EOF) & 0 & 9 & 17 \\
Tỉ lệ lỗi HTTP       & 0\%   & \warn{5,0\%} & \bad{13,0\%} \\
\bottomrule
\end{tabular}
\caption{So sánh kết quả tổng quan theo mức tải.}
\end{table}

\begin{table}[h]
\centering
\renewcommand{\arraystretch}{1.25}
\begin{tabular}{@{}l c c c@{}}
\toprule
\rowcolor{head!12}
\textbf{Thời gian phản hồi AI} & \textbf{Baseline} & \textbf{30 user} & \textbf{50 user} \\
\midrule
Nhanh nhất (min)   & --    & 0,59 s  & 1,34 s \\
Trung vị (median)  & 2,5 s & 19,43 s & 23,26 s \\
Trung bình (avg)   & 2,5 s & 20,90 s & \bad{36,28 s} \\
p90                & --    & 32,71 s & \bad{108,40 s} \\
p95                & --    & 55,11 s & \bad{139,16 s} \\
Chậm nhất (max)    & 2,5 s & 89,28 s & \bad{157,20 s} \\
\bottomrule
\end{tabular}
\caption{Phân bố thời gian phản hồi AI theo mức tải (đơn vị: giây).}
\end{table}

\section{Phân tích}

\subsection{Thông lượng đạt trần}
Thông lượng gần như không đổi giữa 30 và 50 người dùng ($0{,}75$ so với $0{,}70$ req/giây).
Điều này cho thấy hệ thống đã \textbf{đạt trần xử lý} ở mức $\approx 0{,}7$--$0{,}75$ request/giây.
Việc tăng thêm người dùng \emph{không} làm tăng thông lượng mà chỉ kéo dài hàng đợi,
khiến độ trễ và tỉ lệ lỗi tăng vọt.

\subsection{Nút thắt cổ chai: mô hình ngôn ngữ cục bộ}
Nguyên nhân chính là \textbf{Ollama chạy cục bộ xử lý gần như tuần tự}. Với 2 lần gọi LLM
cho mỗi câu hỏi, khi có 30--50 người dùng cùng lúc, các yêu cầu xếp hàng dài: thời gian phản hồi
tăng từ $2{,}5$ giây (1 người) lên trung vị $\approx 23$ giây và p95 tới $139$ giây (50 người)
--- tức chậm gấp \textbf{hơn 50 lần}.

\subsection{Lỗi rớt kết nối}
Các lỗi đều thuộc dạng \texttt{EOF} (máy chủ đóng kết nối khi quá tải). Tỉ lệ lỗi tăng từ
$5\%$ (30 người) lên $13\%$ (50 người), do kết hợp: hàng đợi Ollama, giới hạn kết nối tới Neon,
và Uvicorn chỉ chạy \textbf{một worker} duy nhất.

\section{Kết luận}
\begin{center}
\fbox{\parbox{0.92\textwidth}{\centering
\textbf{Điểm bão hoà của hệ thống hiện tại $\approx$ 5--10 người dùng đồng thời.}\\[3pt]
Dưới 10 người dùng: phản hồi nhanh, gần như không lỗi.
Từ 15 người dùng: độ trễ tăng phi tuyến, bắt đầu rớt kết nối.
Ở 30--50 người dùng: p95 đạt 55--139 giây và 5--13\% request thất bại ---
\textbf{không đạt trải nghiệm chấp nhận được} cho người dùng thực.}}
\end{center}

\section{Khuyến nghị}
\begin{longtable}{@{}>{\bfseries}c p{9.5cm} c@{}}
\toprule
Ưu tiên & Giải pháp & Kỳ vọng \\
\midrule
\endhead
\bad{Cao} & Chạy Uvicorn nhiều worker (\texttt{-{}-workers 4}) hoặc dùng Gunicorn & Giảm rớt kết nối \\
\bad{Cao} & Tăng song song Ollama (\texttt{OLLAMA\_NUM\_PARALLEL}) hoặc chuyển sang API cloud (Anthropic/OpenAI) cho tải cao & Gỡ nút thắt LLM \\
\warn{TB} & Cache câu hỏi và câu lệnh SQL lặp lại & Giảm số lần gọi LLM \\
\warn{TB} & Cấu hình connection pool tới Neon hợp lý, tăng timeout & Giảm lỗi EOF \\
\ok{Thấp} & Thêm hàng đợi và thông báo ``đang xử lý'' cho người dùng & Cải thiện trải nghiệm khi tải cao \\
\bottomrule
\end{longtable}

\vfill
\noindent\rule{\textwidth}{0.4pt}\\
\small Báo cáo được lập dựa trên số liệu thực đo từ k6 (\texttt{k6\_summary.json},
\texttt{k6\_summary\_50.json}). Kịch bản kiểm thử: \texttt{loadtest/ai\_load\_test.js}.

\end{document}
